\section{Measurement noise}

The timing results come from a single machine running under WSL2, and I decided
early not to pretend otherwise. There is no control over the CPU frequency
governor under WSL, so the samples carry a few milliseconds of jitter that no
amount of averaging removes. The kernels run in tens of milliseconds, so for the
kernels whose structural change is small, the runtime difference genuinely falls
inside the measurement spread.

The tempting move was to inflate the workloads until every kernel showed a clean
runtime difference. I did not do that, because it would manufacture a runtime
story that the passes did not earn. Instead the harness reports the median with
the interquartile range, the runtime figure draws that range as an error bar so a
reader can see which deltas are signal, and the prose labels the noise dominated
kernels as structural only. One kernel, the arithmetic heavy one with real
redundancy, shows a runtime improvement clearly outside the spread, and that is
the one case where a runtime claim is warranted.

The timing hygiene I could apply, I did: a fixed repetition count with warmup
runs discarded, the median rather than the mean so a single slow sample does not
dominate, and pinning to a fixed set of cores with \texttt{taskset}. The manifest
records the repetition count, the core list, and a note that the governor is
unavailable, so the conditions travel with the numbers. Better numbers would come
from bare metal Linux with a pinned governor, which is noted as future work
rather than quietly assumed.
